\chapter{Observation processes and automatic organization of banking data}
\label{r3:observation}
\section{What the bank records determines what can be learned}
A causal estimator sees records, while the intervention acts on customers,
contracts and markets. Let $X^*$ denote the scientific variables, $R$ their
observation indicators, $L$ an entity-linkage assignment, and $O$ the stored
records. An observation kernel $K$ maps the full law into the recorded law:
\begin{equation}
 p_O(o)=\int K(o\mid x^*,r,l)\,p(x^*,r,l)\,d(x^*,r,l).
 \label{r3:observation-kernel}
\end{equation}
The target $\theta$ may depend on interventions in the law of $X^*$,
not merely on the distribution of $O$. Identification therefore asks
whether all admissible full models with the same $p_O$ give the same
$\theta$. A cleaning transformation $T(O)$ cannot distinguish full models
that already agree on $O$: their pushforward laws under $T$ also agree.
This elementary argument explains why a more complete-looking table
cannot by itself resolve missingness, measurement error or causal ambiguity.

Useful automation instead makes the observation assumptions explicit,
checks their observable implications and preserves information for an
appropriate estimator. The aim is to automate reproducible construction
of an analysis view while retaining a scientific record of what that view
means. Data lineage proves how a variable was computed; it does not prove
that the variable blocks a confounding path.

\section{Event, availability and revision clocks}
\label{r3:availability}
For a decision at time $t$, let $\mathcal F^{\rm bank}_{t-}$ be the
information actually available immediately before assignment. A deployable
policy $g_t$ and its inputs must be measurable with respect to that
information. This is different from having an event timestamp earlier
than $t$. A January income field entered in March may describe January
but was unavailable to a February decision. A revised arrears balance
may replace an older value in today's warehouse while the historical
policy acted on the older version.

Represent a record by its entity, event time, availability time, validity
interval, revision identity, value, unit and source. An as-of view at
$t$ first restricts records to availability before $t$, then chooses
the applicable version for the scientific event interval. For example,
a deposit balance dated 31 January, posted on 2 February and corrected
on 10 February supplies no input to a 1 February offer; it supplies
its initial value on 5 February and its corrected value on 12 February.
Whether the outcome should use the eventually reconciled balance or
the historical published balance is a separate estimand choice.

There are also two scientifically different uses of retrospectively
measured variables. An eventual audit can validly measure a truly
predecision confounder for an explanatory analysis under its measurement
model, even if that measurement was unavailable to the historical policy.
It cannot then be used as a prospective policy feature. Conversely,
a postdecision credit score that incorporates the offered credit limit
is a descendant of treatment; its event label does not make it a valid
baseline adjustment variable. The product must preserve both causal
time and available information.

For longitudinal evaluation, rolling outcomes extend beyond their origin.
If $Y_{t,h}$ aggregates cash flows through $t+h$, training rows with
earlier origins can still use outcomes that occur after an evaluation
decision. A prospective split needs outcome availability as well as
feature availability. An offline estimator of a stationary law may use
a different split justified by dependence assumptions, but it answers
a different validation question.

\section{Missing outcomes: identification before imputation}
\label{r3:missingness}
First isolate missingness from treatment assignment. Suppose $X$ is fully
observed, $Y$ is a well-defined outcome for every unit, and $R=1$ means
that $Y$ is observed. Define $\pi(x)=P(R=1\mid X=x)$. Under
$Y\perp R\mid X$ and $\pi(X)>0$ on the target support,
\begin{equation}
 \mu=\E[Y]=\E[m(X)]=\E\!\left[\frac{RY}{\pi(X)}\right],
 \quad m(x)=\E[Y\mid X=x,R=1].
 \label{r3:missing-mean}
\end{equation}
Condition on $X$: the weighted observed mean is
$\pi(X)\E[Y\mid X,R=1]/\pi(X)=\E[Y\mid X]$.
For candidates $\widetilde m,\widetilde\pi$, augmentation gives
\[
 \widetilde\mu=\E\!\left[
 \widetilde m(X)+\frac{R}{\widetilde\pi(X)}\{Y-\widetilde m(X)\}\right],
 \qquad
 \widetilde\mu-\mu=
 \E\!\left[\frac{\widetilde\pi-\pi}{\widetilde\pi}
                              (\widetilde m-m)\right].
\]
This project derivation has the same observation-score structure as
Chapter~\ref{r3:learning}. It does not prove missing at random.

When treatment $A$ and covariates $X$ are observed for everyone, static
exchangeability for $A$ plus $Y\perp R\mid A,X$ identifies
$\E[Y(a)]$ through $\E_X\E[Y\mid A=a,X,R=1]$; the corresponding
weight is $\mathbf1(A=a)R/\{e_a(X)\pi(a,X)\}$.
If the bank observes outcomes only while accounts remain open,
sequential censoring replaces this one-step model and requires the
history-specific assumptions in Chapter~\ref{r3:longitudinal}.
If an outcome is not defined for rejected applications without
specifying loan origination, treating it as an ordinary missing value
silently introduces a different intervention.

\subsection{A missingness law that no imputer can learn from these records}
Suppose $P(R=1)=1/2$ and observed binary outcomes are evenly split.
One full model has $P(Y=1\mid R=0)=0$ and mean $1/4$; another has
$P(Y=1\mid R=0)=1$ and mean $3/4$. They have exactly the same
recorded law $(R,RY)$. The general bounded-outcome identified interval is
\[
 P(R=1)\E[Y\mid R=1]
 \ \le\ \E[Y]\ \le\
 P(R=1)\E[Y\mid R=1]+P(R=0),\qquad 0\le Y\le1.
\]
Every value between the endpoints can be attained by a suitable
unobserved outcome distribution. Filling the gaps with the observed
mean selects $1/2$ through an assumption; prediction validation on
observed cases cannot distinguish these models. A sensitivity parameter
$\delta(x)=\E[Y\mid X=x,R=0]-\E[Y\mid X=x,R=1]$ instead yields
the explicit correction $\E[(1-\pi(X))\delta(X)]$ to the MAR mean.
Its admissible range can be connected to bounded outcomes or external
follow-up data.

Missingness also changes graph discovery. Test-wise complete cases for
different triples need not come from one common observational law;
conditioning on observation can open collider paths.
\citet[Section 3, Assumptions 1--4]{tu2019} studies this problem under
causal sufficiency, faithfulness/observability restrictions and specific
constraints on missingness indicators, including no self-masking.
Its correction procedure is not a general MNAR solution. If customers
with the greatest unrecorded distress avoid reporting their distress,
the no-self-masking premise is precisely the questionable restriction.

\section{Measurement error is another inverse problem}
\label{r3:measurement}
Consider binary true receipt $A$ and recorded receipt $W$. Under a
known nondifferential error model,
$P(W=w\mid A=a,Y,X)=K_{wa}(X)$. For each outcome value and covariate
stratum, the observed joint probabilities obey
\begin{equation}
 \begin{pmatrix}p(W=0,Y\mid X)\\p(W=1,Y\mid X)\end{pmatrix}
 =K(X)\begin{pmatrix}p(A=0,Y\mid X)\\p(A=1,Y\mid X)\end{pmatrix}.
 \label{r3:misclassification}
\end{equation}
If false-positive and false-negative probabilities are $f,g$, then
$K=\left(\begin{smallmatrix}1-f&g\\ f&1-g\end{smallmatrix}\right)$
has determinant $1-f-g$. Known invertible $K$ identifies the true
joint law algebraically; a nearly singular $K$ amplifies error.
Unknown $K$ is an additional identification problem requiring validation
data, repeated measurements or other restrictions. Correcting recorded
receipt does not by itself remove confounding of true receipt.

For an exact example, let true receipt be randomized with probability
$1/2$, and let $P(Y=1\mid A=0)=1/5$ and
$P(Y=1\mid A=1)=4/5$. Its causal contrast is $3/5$.
With symmetric misclassification $f=g=1/5$, the recorded conditional
means are $8/25$ and $17/25$, so the naive contrast is $9/25$.
Applying the known matrix inverse recovers $3/5$. If $f=g=1/2$,
$W$ contains no information about $A$ and the inverse does not exist.
This is a finite version of the operator/rank distinction in
Chapter~\ref{r3:inverse-learning}; more rows cannot restore lost rank.

High-dimensional denoising can use additional structure.
\citet[Sections 3.2--3.3 and 4]{corrupt2022} studies corrupted covariates
under approximate low-rank/repeated-measurement structure, with a
specified corruption model and inference that accounts for cleaning.
The main model uses covariates plus additive noise and masking; it does
not cover arbitrary deterministic rounding or arbitrary informative
missingness. The scientific contribution is a joint restriction on
cleaning, target estimation and error rates. Applying PCA to any bank
table and then using an ordinary standard error is not that procedure.

\section{Entity uncertainty and the unit of intervention}
\label{r3:entity-uncertainty}
A household may own several customer identifiers, an SME may merge,
and a branch portfolio may transfer to another legal entity. The unit
for assignment can therefore differ from the unit for reporting and
the unit over which interference occurs. A household-level promotion
analyzed as independent accounts can both misstate treatment versions
and exaggerate sample information. A stable identifier is not evidence
of a stable economic unit across a merger.

Let $L$ denote a candidate linkage and $\theta_L$ the effect estimate
or parameter under it. If a justified joint probabilistic model supplies
linkage weights, its posterior uncertainty satisfies
\begin{equation}
 \Var(\theta\mid O)=
 \E[\Var(\theta\mid O,L)\mid O]
 +\Var(\E[\theta\mid O,L]\mid O).
 \label{r3:linkage-variance}
\end{equation}
The second term disappears when a most-likely linkage is treated as
certain. This is the law of total variance, not a guarantee that an
arbitrary matcher's confidence scores are calibrated posterior weights.
For example, two equally weighted linkage models with effect means
$0$ and $2$ and conditional variance $1/4$ give total variance $5/4$.
Reporting $1/4$ misses the between-linkage component.

An alternative is sensitivity across a defensible set of linkages,
returning a union of their identified/estimated effect sets. This does
not require invented posterior probabilities but can be wide. Averaging
linked rows first generally differs from averaging the resulting causal
functional because propensity scores, conditional laws and target
weights are nonlinear. The entity service should preserve uncertain
links, historical validity and their effect on the target population.

\section{Which transformations preserve the scientific question?}
\label{r3:data-transformations}
A one-to-one covariate transformation preserves its sigma-field:
$\sigma(f(X))=\sigma(X)$ when the inverse is measurable. Conditional
exchangeability and conditional expectations can then be translated
without losing information. A coarsening generally cannot. A special
exception is a balancing score under the assumptions derived in
Section~\ref{r3:representation}; arbitrary low-dimensional embeddings
do not enjoy that property.

Outcome transformations also change the target. An effect on
$\E[\log Y]$ is not the log of an effect on $\E[Y]$.
Portfolio sums, account averages and balance-weighted averages assign
different units and weights. In finance, a useful reward must reconcile
interest margin, fees, credit losses, funding cost, servicing cost and
capital charges over the same intervention horizon. Internal transfers
can be a gain to one branch and no gain to the banking group. Accounting
identities should be retained as constraints, not supplied as independent
variables to a faithfulness-based discovery procedure.

For example, if ending deposits equal starting deposits plus net flows,
then the three variables are deterministically related. Their covariance
is singular and their independences need not behave like a generic
noisy DAG model. Choose primitive quantities or an explicit constrained
measurement model. Likewise, mechanically constructing an interest
income variable as balance times rate does not establish either economic
causal direction from a discovered association.

At large scale, an immutable analysis recipe can specify event/availability
filters, revision policy, entity map, action history, eligibility, outcome
windows, units and observation indicators. A generated query should be
checked for future information, join multiplicity, unsupported follow-up,
unexpected missingness and reconciliation to source totals. Partitioning
by entity/date and streaming score reductions can preserve this recipe;
downsampling, aggressive screening and global imputation can change the
scientific information. Their consequences must be tested separately
from query execution time.

\section{A finite candidate search with explicit abstention}
\label{s4:organization-algorithm}
An initial automatic organization operation can be specified without
claiming to infer arbitrary economic meaning. Let a candidate observation
plan be $c=(k,e,a,r,g,u)$: record key $k$, event and availability clocks
$e,a$, revision relation $r$, aggregation $g$, and entity map/unit $u$.
The input includes a finite grammar $\mathcal C$, raw versioned records,
decision identifiers and times, declared units, and constraints with their
provenance. The grammar is itself part of the scientific input. A method
cannot recover an interpretation that the grammar excludes.

The implemented reference grammar has two timestamp assignments, two
record-key hypotheses (source/record or source/entity/record), latest-version
versus append semantics, stock versus flow aggregation, and account versus
a supplied candidate household map. This gives 32 initial candidates.
The engine tests key uniqueness, finite values, currency compatibility,
declared column roles, and any supplied chronology, revision-publication
or accounting constraints. For example, $e\le a$ can distinguish two
timestamp assignments only if the analyst has declared that this feed
records realized events before their publication. Scheduled future events
do not obey that restriction. An external closing-minus-opening ledger
identity can distinguish flow summation from a last stock only if that
identity and its independent measurement are actually supplied.

For each surviving candidate and each requested decision time $t$, first
keep versions with $a<t$. Within its record identity, select the largest
eligible version; then restrict its economic date to the requested
half-open interval $[b,t)$. A correction moving an event outside the
interval must not resurrect the older version. Distinct decision identifiers
and distinct events remain distinct even when their amounts match.
A flow is the sum of distinct eligible increments. A stock is its latest
eligible economic observation; component accounts at that instant can be
summed only under the supplied additive household interpretation. Missing
support returns a missing value, not a fabricated zero. The output retains
the source/entity/record/version identifiers of every selected term.

Let $V_c$ be this materialized view including units and lineage. Hard
constraints define $\mathcal C_0\subseteq\mathcal C$. A transparent
inspection score is
\[
 s(c)=\sum_j\ind\{\text{candidate field }j\text{ agrees with supplied hint }j\}.
\]
Hints order inspection; they cannot eliminate a surviving material
alternative. Partition the survivors by equality of requested values,
lineage, aggregation meaning and economic unit. Return incompatible if
none survive, accept the current requested view if one class remains,
and abstain from a unique view otherwise. Even accepted equality for the
current decisions does not identify a rule for unseen data. Changing the
request can separate previously equivalent plans.

This construction has a simple conditional correctness statement. If the
true plan belongs to the grammar and every hard constraint is valid for
it, that plan survives. If all surviving plans agree on the required view,
the accepted view equals its oracle counterpart. No probabilistic calibration
is needed for this implication. Incorrect metadata or an incomplete grammar
invalidates its premise. A large hint score cannot estimate the probability
that those premises are true. Under ambiguity, fit the scientific functional
separately in each view and retain the resulting sensitivity set; do not
average the tables or treat hint scores as linkage probabilities.

\subsection{An evaluation that reaches the effect estimate}
The synthetic experiment fixes the oracle before constructing candidates.
For 240 customers draw independent standard-normal $X,W$, assign an offer
with probability $\operatorname{expit}(.7X)$, and generate
$Y=20A+30X+\epsilon$, $\epsilon\sim N(0,25)$. The predecision record
contains $X$; a later revision contains $X+2A$ and is unavailable at the
decision. OLS on $(1,A,X)$ estimates the constant effect of 20 in this
specified law. Using the eventual revision instead gives the population
equation $Y=-40A+30(X+2A)+\epsilon$: the nominal cleaning choice changes
the adjustment variable into a treatment descendant.

Across 300 repetitions, labeled revisions and the valid chronology case
accept the oracle view in every run, with effect RMSE about $0.68$--$0.69$.
Retrospective revision adjustment has RMSE about $60$. In the unlabeled
aggregation case an earlier independent record $W$ leaves last-stock $X$
and flow-sum $X+W$ as materially different possibilities. In the linkage
case pairing two accounts changes the available confounder to their sum.
Both cases abstain in every run and retain the oracle. Forcing the top
hint-ranked candidate gives RMSE $10.19$ and $10.59$, while the mean
spread between alternative estimates is $9.77$ and $10.13$. Those spreads
are observation sensitivity, not confidence intervals.

An accounting case has increments $W$ and $X-W$, with an independently
declared reconciled total $X$. The constraint selects the flow plan in every
run, with RMSE $0.70$. The additional total provides the distinguishing
information; the algorithm did not discover it from column names.
Zero wrong acceptances among these fixtures is conditional evidence for
this grammar and generator. It does not establish population error rates
for previously unseen bank feeds. With zero errors in 300 trials, the
two-sided 95\% Wilson upper bound is about $1.27\%$ for the particular
repeated experiment, not zero.

The bounded search costs $O(|\mathcal C|RD)$ in a direct implementation
with $R$ records and $D$ decisions. Indexing by entity and availability
time, caching common candidate projections and expressing eligible-version
selection in a columnar engine reduce repeated scans without changing the
semantics. Probabilistic record linkage is a separate extension requiring
a calibrated joint linkage/outcome model or an explicitly conservative
linkage set. The practical stopping point is a reproducible observation
plan or a short list of decision-relevant alternatives that an analyst can
resolve with additional metadata or measurement.


\section{Automation with a defensible stopping point}
The data layer can propose schemas, map candidate variables, construct
as-of joins, expose revisions and generate reproducible views. It can
identify an ambiguous action timestamp or a proxy measured after the
decision and explain which analysis is affected. It cannot establish
exclusion, exchangeability, missing-at-random or a unique household
definition solely from column names. Those inputs belong to the
scientific model and may remain alternative hypotheses.

The output is consequently more than a clean matrix: it is a defined
population and observation law, an analysis view with provenance, and
explicit unresolved measurement or linkage alternatives. Passing that
object to discovery, longitudinal estimation or state-space fitting
preserves the connection between the bank's intervention and the
numbers ultimately reported.

\section{Generating observation hypotheses and asking useful questions}
\label{s5:observation-generation}
The preceding finite grammar supplies an exact conditional reconstruction
operation. A reusable product must also explain where candidates come from.
The new profiler extracts datetime columns as possible clocks, string
columns as possible identifiers, small nonnegative integer fields as possible
versions, and finite floating fields as possible values. Candidate entity
fields must contain the supplied decision-unit identifiers. The search forms
record keys of up to two additional fields, pairs event and availability
clocks, and combines version, append/latest and stock/flow alternatives.
These are type and membership hypotheses; economic meaning is not inferred
from the arbitrary column names.

For decision cutoff $c$, a latest-version plan first restricts records to
availability strictly before $c$, then selects the last eligible revision
within its record identity, then applies the event-time window. This order
matters when a correction changes the event timestamp. A stock plan selects
the latest event instant and rejects multiple conflicting stock values;
a flow plan sums its eligible records. Chronology
$t_{\rm event}\le t_{\rm available}$ is imposed only for a feed explicitly
declared to contain realized events. A scheduled future contract can validly
be known earlier, so the same inequality is not a universal cleaning rule.
Unique revision identity, supplied semantic facts and supported search bounds
determine rejection. Empty survivors mean incompatibility within this
candidate language; they do not certify that no correct interpretation exists.

Let $\mathcal C$ be surviving plans and $\theta_c$ the downstream diagnostic
effect computed under plan $c$. A semantic question $q$ partitions candidates
by possible answer $a$. The implemented choice minimizes
\[
 \max_a\{\max_{c\in\mathcal C_{q,a}}\theta_c-
                 \min_{c\in\mathcal C_{q,a}}\theta_c\},
\]
then the worst remaining number of material views. With no justified
probability over interpretations this is a minimax ambiguity criterion,
not expected information gain. A trusted answer filters the candidates;
the algorithm never supplies its own answer from the ranking. Using outcomes
to prioritize questions is exploratory. If the resulting choice changes
the final statistical analysis, inference must account for that selection
or use independent metadata/pilot information. Small effect spread today
does not prove semantic equivalence under future decisions or populations.

\subsection{Conditional correctness and missing interpretations}
Suppose a correct plan $c^\star$ is generated, all hard constraints and
analyst answers are true of it, and accepted plans agree on the requested
material observation. Then $c^\star$ survives every filtering step and the
accepted observation equals its projection. This induction proves the
conditional correctness claim. It does not prove candidate inclusion,
answer truth or population-level semantic recovery. Consequently the workflow
records proposed hypotheses, supplied facts, rejected reasons and unresolved
views independently. A contradiction triggers a new candidate/search request;
it must not silently promote the highest-ranked survivor.

The test fixtures rename and reorder every field, include future revisions,
and vary stock/flow interpretation and alternative value roles. Across 80
repetitions in each supported case, the oracle observation survives and is
resolved without wrong acceptance. One oracle-supplied answer resolves stock
and flow cases; two resolve the value-role ambiguity. An intentionally
unrepresented three-field record key yields incompatibility in every run,
with no false acceptance. These are synthetic heterogeneous schemas with
trusted fixture answers, not measured analyst accuracy or burden in a bank.
The downstream adjusted-effect spread determines question priority, while
the answer's provenance remains the actual source of semantic resolution.

\section{Linkage as a measurement model with audit uncertainty}
\label{s5:linkage}
A resolved table schema does not guarantee the correct customer is linked.
For a transparent randomized example let $A\sim\operatorname{Bernoulli}(1/2)$,
$Y=\tau A+\epsilon$, and let $\widetilde Y=(1-M)Y+MY'$ with
$M\sim\operatorname{Bernoulli}(\alpha)$ independent of $(A,Y,Y')$.
The replacement $Y'$ is independent of the current treatment and has the
same marginal outcome distribution. The known-randomization score is
$Z=2(2A-1)\widetilde Y$. Then
\[
 \E[Z]=(1-\alpha)\tau,\qquad
 \hat\tau=\bar Z/(1-\hat\alpha).
\]
The replacement contribution is zero because
$\E[2(2A-1)Y']=0$. With an independent truthful audit of $M_v$ links,
the first-order expansion is
\[
 \hat\tau-\tau=\frac{\bar Z-(1-\alpha)\tau}{1-\alpha}
 +\frac{\tau}{1-\alpha}(\hat\alpha-\alpha)+o_p(N^{-1/2}+M_v^{-1/2}),
\]
and its variance is
\begin{equation}
 \frac{\Var(Z)}{N(1-\alpha)^2}
 +\frac{\tau^2\alpha(1-\alpha)}{M_v(1-\alpha)^2}.
 \label{s5:link-variance}
\end{equation}
Plugging $\bar Z$ into the second term gives the implemented
$\bar Z^2\hat\alpha(1-\hat\alpha)/\{M_v(1-\hat\alpha)^4\}$.
An overlapping audit would add cross covariance. Near $\alpha=1$, the
denominator is weak and a normal ratio interval is inappropriate. This
special attenuation calculation is a project construction; the broader
linked-data causal analysis requires conditional mismatch and independent
truthful-audit restrictions \citep[Assumptions A1--A5]{slawski2025}.
A generic fuzzy-match confidence score is not an estimate of this $\alpha$.

With $N=3,000$, $M_v=80$, $\alpha=.2$ and true effect eight, 400 repetitions
give naive bias $-1.5923$ and zero nominal 95\% coverage. Joint audit
correction gives bias $.0337$ and .955 coverage. Treating the estimated audit
rate as fixed gives .615 coverage. A biased audit drawn from links with
mismatch rate .05 gives bias $-1.2519$ and .015 coverage despite retaining
its sampling variance. Propagating uncertainty is necessary; a biased
measurement model remains biased after propagation.

For richer entity resolution, preserve a distribution or set of plausible
link assignments together with its calibration source. Averaging causal
estimates over match weights is justified only by a coherent joint linkage
and outcome/assignment model; conditioning on a match chosen using outcome
information can introduce selection. The product should offer either a
qualified audited correction, sensitivity over linkage laws, or separate
candidate analyses with unresolved status. The synthetic audit calculation
shows why exporting a single hard customer ID and discarding the matching
evidence is not an adequate scientific interface.


\section{Conflicting source meaning and calibrated candidate histories}
\label{s6:linked-histories}
Consider a deposit feed with business date, warehouse arrival date and
revision number, plus a separate account/customer ownership table. A row
for business day 10 may arrive on day 20, after the day-15 decision. A
business-date filter then exposes a revision unavailable to the decision
maker. Joint ownership creates another ambiguity: copying an account
balance to both owners changes the economic unit and can double exposure.
These are observation-model questions before any causal estimator is called.

Five source-owner fixtures are specified separately from the interpretation
code. A current trusted arrival-clock statement excludes the late revision
and gives balance 90. Conflicting trusted arrival/business-clock statements
retain balances 90 and 110; missing clock documentation does the same. A
customer request with two joint owners and no allocation rule abstains.
Finally, a plausible but incorrect accepted migration specification yields
the unique wrong balance 110. The algorithm reports conditional acceptance
in that case, and the external fixture evaluation records wrong acceptance.
The result is not hidden by defining every supplied answer as correct.

Verified provenance means that the stated owner supplied the assertion, not
that the assertion's economic meaning is true. Schema profiles can generate
candidate clocks and entity relations; they cannot establish first
availability or ownership allocation by themselves. An analyst question
must therefore request the missing semantic fact and identify its source.
If pilot outcomes rank useful questions, the ranking and answers are frozen
before final evaluation. The executed final linkage experiment below uses
predeclared strata and independent audit information, with no final-outcome
question selection.

\subsection{Whole candidate maps and a recoverable mixture}
A scalar replacement rate is inadequate when plausible candidate entities
have different responses. Use two candidate portfolio histories at each
origin, with the entire response vector assigned consistently under each
map. In linkage stratum $s$, the two histories are correctly ordered with
probability $p_s$ and swapped otherwise. This is a one-to-one pair assignment,
not independent per-column linking. The match event is independent of the
complete potential histories conditional on the declared stratum and
matching information. That strong noninformative-linkage restriction is
part of identification; a name-match score is not evidence for it.
The executed map is redrawn by origin, while retaining all horizons within
that origin together. It models an independently reconstructed candidate
window, not a persistent customer alias across the entire calendar. A fixed
erroneous crosswalk or time-dependent matching mechanism would need its own
joint map law and covariance; the following evidence does not qualify them.

Let $\theta_{sA},\theta_{sB}$ be the true response vectors of the two
portfolios to the same price innovation. The two observable candidate
response vectors satisfy
\begin{equation}
 \begin{pmatrix}\eta_{s1}\\\eta_{s2}\end{pmatrix}
 =\begin{pmatrix}p_s&1-p_s\\1-p_s&p_s\end{pmatrix}
   \begin{pmatrix}\theta_{sA}\\\theta_{sB}\end{pmatrix},\qquad
 \theta_{sA}=\frac{p_s\eta_{s1}-(1-p_s)\eta_{s2}}{2p_s-1}.
 \label{s6:pair-mixing}
\end{equation}
This corrects two different candidate responses rather than assuming every
wrong entity has zero effect. An independently audited set of pair matches
estimates each $p_s$. At $p_s=1/2$ the matrix has rank one: these mean
equations identify only the sum of the two responses. In the symmetric
pair-swap observation law, exchanging portfolio labels preserves the
observed law; more unlabelled histories cannot resolve that named-portfolio
ambiguity. Higher mixture moments under additional restrictions could still
identify more than these mean equations do. If probabilities are not
calibrated, plausible maps produce a set of targets, not a probability-
weighted point estimate justified by convenience.

The funding experiment fits candidate conditional means separately by
stratum using an independent pilot and the prespecified driver basis.
The potentially misassigned lagged balance is excluded from this basis.
Under the stationary Gaussian driver law, conditional means of every past
driver and its centered square given the current driver are linear and
quadratic, respectively. Summing the stable balance recursion therefore
keeps the candidate outcome means in the driver basis after that lagged
balance is marginalized. This justification depends on the stated law;
it does not license dropping a confounding history in a new bank dataset.
Treatment residuals, candidate outcomes and stratum indicators produce
the candidate response moments. The conditional mean of treatment does
not depend on stratum, but allowing stratum interactions in its learner
preserves the common nuisance interface. Estimating candidate outcomes with
one pooled stratum mean could destroy the orthogonality needed for
stratum-specific inference; the executed design retains those interactions.

\subsection{Uncertainty from dates and calibration}
For fixed target weights $w_s$, let
$\theta=\sum_sw_s\theta_{sA}$. Stack all candidate response influences by
calendar date; the same dates, horizons and candidate histories induce
cross-component covariance $V_\eta$. With independent audits of fixed sizes
$M_s$, delta differentiation gives
\begin{align}
 d\theta_{sA}
 &=\frac{p_s}{2p_s-1}\,d\eta_{s1}
   -\frac{1-p_s}{2p_s-1}\,d\eta_{s2}
   +\frac{\eta_{s2}-\eta_{s1}}{(2p_s-1)^2}\,dp_s,\nonumber\\
 V_\theta
 &=J_\eta V_\eta J_\eta'
   +\sum_s j_sj_s'\frac{p_s(1-p_s)}{M_s},\qquad
 j_s=w_s\frac{\eta_{s2}-\eta_{s1}}{(2p_s-1)^2}.
 \label{s6:linkage-variance}
\end{align}
The calibration term is a matrix outer product: a single map probability
affects the entire response curve. Both $T$ and $M_s$ must grow in the
regular joint-sample approximation; a fixed small audit is not washed out
by arbitrarily many production records. Bias in the audit changes the
inverse map itself. Covariance correction cannot remove that bias.
If the Wilson audit interval contains $1/2$, the regular reference withholds the
regular inverse-map interval. This is an explicit availability rule, not
a claim of uniform post-selection coverage.
Section~\ref{s61:linkage-sets} supplies a separate projection procedure that
retains unbounded and disconnected financial confidence sets near rank
failure. It does not extend this delta approximation by changing its label.

In the registered experiment $p_0=.95$, $p_1=.8$, the true second-portfolio
price coefficient is .5 rather than the focal portfolio's 2, and both
histories share common shocks. With 720 evaluation dates and 100 audited
pairs per stratum, all 400 repetitions are available and financial coverage
is .9625; at 400 audits it is .950. Treating calibration as fixed gives
.960 and .950 respectively. Thus calibration is propagated, but its measured
variance contribution is modest relative to the calendar noise in this
design. The naive first-candidate analysis has bias about $-68{,}000$ and
$-80{,}000$ currency units in the two runs, with coverage .900 and .9075.
The correctly mapped financial target is about $375{,}149$ under the fixed
accounting inputs. These wide intervals are evidence about a difficult
dynamic observation problem, not a promised precise business case.

This construction is a bounded project derivation. Audit-based mismatch
correction provides one related approach \citep{slawski2025}; joint
probabilistic linkage and causal modeling offers a different route
\citep{guha2022} when a
well-specified joint likelihood and linking covariates are defensible. Such
a joint model may let outcomes inform linkage. Its posterior target and
feedback assumptions differ from this independent-audit frequentist
experiment; posterior match draws cannot simply be treated as additional
independent customers. The product records the selected mapping law, audit
population, covariance and unresolved semantic alternatives together.


The publisher version of \citet[Section 2.3, equation (2.4)]{guha2022}
models a constrained one-to-one map $z$, agreement fields and outcome
regressions jointly. The linked-record likelihood includes
$f_1(y_i\mid x_j,a_j)$; unlinked outcomes have $f_2(y_i)$. Thus the
posterior odds of a map contain the outcome likelihood ratio as well as
field agreement. Section 2.1 requires strong treatment ignorability and
positivity; the model cannot repair hidden action confounding by better
linkage. The journal version explicitly separates linking variables from
causal covariates and discusses conditioning/altering the likelihood when
they overlap. Its linked-population treatment effect and within-file
no-duplicates assumption differ from a repeated banking history with
joint ownership. These differences are why a production adapter must
represent the map law and target population, rather than passing arbitrary
posterior links into an otherwise unchanged temporal estimator.

\section{Weak linkage identification is a financial inverse problem}
\label{s61:linkage-sets}
The inability to distinguish two candidate histories is structurally related
to a weak instrument. In Equation~\eqref{s6:pair-mixing}, the eigenvalues of
the assignment matrix are $1$ and $r_s=2p_s-1$. The portfolio sum survives
even when the difference is unidentified. A Wilson interval containing
$1/2$ usefully stops the regular inverse, but that screening rule does not
construct an identification-robust confidence set. Nor does a large number
of production dates create audit information. The relevant analogy with
Section~\ref{r3:weak-iv} is to build uncertainty for observable quantities
before dividing by an uncertain identification coefficient
\citep[Section 2D, Theorem 3]{staiger1994}.

Use the fixed financial coefficient vector $k$ from
Equation~\eqref{s6:funding-account}, including the cost of paying the higher
first-month rate on induced balances. Within stratum $s$, define
\[
 m_{s1}=k'\eta_{s1},\quad m_{s2}=k'\eta_{s2},\quad
 S_s=\frac{m_{s1}+m_{s2}}2,\quad D_s=\frac{m_{s1}-m_{s2}}2.
\]
Let $v_{sA}=k'\theta_{sA}$ and $v_{sB}=k'\theta_{sB}$. Adding and
subtracting the two mixing equations gives
\begin{equation}
 S_s=\frac{v_{sA}+v_{sB}}2,\qquad
 D_s=r_s\frac{v_{sA}-v_{sB}}2=r_s(v_{sA}-S_s).
 \label{s61:diagonal-map}
\end{equation}
Thus $v_{sA}=S_s+D_s/r_s$ for $r_s\ne0$. At $r_s=0$, the observable
restriction in these mean equations is $D_s=0$ and $v_{sA}$ is unrestricted
by them because $v_{sB}$ can compensate. The final funding value is
$V=\sum_s w_s v_{sA}-C$, with fixed $w_s=1/2$ and
$C=Nc r B_0$. These are monetary contrasts before the common fixed price
cost is subtracted; no balance stock is reinterpreted as a flow.

This is an inverse problem for response moments. It is not a claim that
every value is compatible with every feature of the full observed mixture
distribution. At $p=1/2$, conditional on an action/history, that distribution
has the symmetric form
$\{P(Y_A,Y_B)+P(Y_B,Y_A)\}/2$. Exchanging the two named portfolio laws
preserves it, so focal identity can remain ambiguous. Additional distributional
restrictions may nevertheless identify unordered components or their
separation and produce a smaller set than the mean equations permit.
The procedure below constructs a conservative outer confidence set without
using those restrictions; it does not compute the sharp identified set of
a fully specified mixture. Its independent audit supplies correctness counts
only. Audited action/outcome histories could directly inform portfolio
responses and would define a richer experiment.

\subsection{A confidence region for regular observable information}
The candidate-response vector is estimated before inverse linkage. Its
calendar influence and HAC covariance are those in
Section~\ref{s6:linked-histories}; every financial reduction uses the entire
within-horizon covariance. Apply the fixed linear transformation above to
obtain $\widehat z=(\widehat S_0,\widehat D_0,
\widehat S_1,\widehat D_1)$ and its covariance $\widehat V_z$.
For total error $\alpha$, allocate $\alpha/2$ to the four observable
coordinates. Marginal bands are
\begin{equation}
 I_j=\left[\widehat z_j-z_{1-\alpha/16}\sqrt{\widehat V_{z,jj}},
                 \widehat z_j+z_{1-\alpha/16}\sqrt{\widehat V_{z,jj}}\right].
 \label{s61:primitive-bands}
\end{equation}
Each uses error $\alpha/8$, giving joint coverage at least
$1-\alpha/2+o(1)$ by Bonferroni when its studentized approximation is valid.
There is no independence assumption between these four estimates.

The approximation concerns candidate means, not a near-singular causal
inverse. It retains the independent pilot, finite correct basis, fixed
horizon, nonvanishing stratum/innovation masses, mixing/moment conditions
and consistent lag covariance established for the funding law. The
bandwidth conditions of \citet[Theorem 2]{newey1986} remain relevant after a
linear transform. Near-unit-root dynamics, unqualified nuisance learning or
a missing evaluation stratum are not repaired by projecting a confidence
region. Positive limiting variances are required for the stated normal
bands; a deterministically known coordinate instead needs an exact
singleton region justified by its model.

For each independent representative audit let
$K_s\sim\operatorname{Binomial}(M_s,p_s)$. Allocate error
$\epsilon=\alpha/4$ to its two-sided probability interval $P_s=[L_s,U_s]$.
For $0<K_s<M_s$, obtain the endpoints from
\begin{equation}
 \begin{split}
 \sum_{j=K_s}^{M_s}\binom{M_s}{j}L_s^j(1-L_s)^{M_s-j}&=\epsilon/2,\\
 \sum_{j=0}^{K_s}\binom{M_s}{j}U_s^j(1-U_s)^{M_s-j}&=\epsilon/2.
 \end{split}
 \label{s61:binomial-inversion}
\end{equation}
Set $L_s=0$ at zero successes, $U_s=1$ at all successes, and
$P_s=[0,1]$ if $M_s=0$. The binomial tail probabilities are monotone
in $p$. For either endpoint the rejecting counts form one tail whose
total probability is bounded by $\epsilon/2$. Taking the union of the
two rejection events proves
$P_p\{p\notin P_s\}\le\epsilon$ for every $p\in[0,1]$.
Discreteness can make coverage conservative. This proof does not require
$M_s$ to grow with production dates; a fixed audit simply leaves a region
that need not shrink. These claims require accurate, noninformatively
sampled audit labels. The audit assumptions in
\citet[Section 2.1, Audit Sample]{slawski2025} make that distinction explicit;
they are not consequences of the interval formula.

\subsection{Projection, coverage and the shape of the result}
Let $\mathcal R$ be the Cartesian product of the four observable bands and
the two audit intervals. For each stratum construct
\begin{align}
 \mathcal V_s&=\{v:\exists S\in I_{S_s},D\in I_{D_s},p\in P_s,
                         \ D=(2p-1)(v-S)\},\nonumber\\
 \mathcal V&=\left(\sum_s w_s\mathcal V_s\right)-C.
 \label{s61:projection}
\end{align}
The sum denotes all sums of one element from each set, not the sum of
arbitrarily selected endpoints. On the event that all primitive quantities
lie in $\mathcal R$, the true $v_{sA}$ satisfy the displayed equations,
including at $p_s=1/2$. Hence the true $V$ lies in $\mathcal V$ on that
same event. The union bound gives
\[
 P\{V\in\mathcal V\}\ \ge\ 1-\alpha+o(1).
\]
This is a project coverage argument, not a new theorem attributed to the
linkage paper or a direct application of the Gaussian Anderson--Rubin test.
The audit region is finite-sample valid; the response region remains
asymptotic. Projection itself needs no lower bound on $|2p_s-1|$.
A uniform claim over a sequence of linkage probabilities additionally
requires uniform validity of the primitive response bands along that
sequence. Pointwise CLTs cannot establish that claim by renaming the result.
The distinction between a finite normal experiment and an asymptotic
approximation is also explicit in
\citet[Section 2.5, equations (9)--(11)]{rambachan2023}.

There is a complete scalar algorithm without a bounded effect grid.
Write $I_S=[a,b]$, $I_D=[c,d]$ and
$R=[2L-1,2U-1]$. On each positive or negative part of $R$,
the extrema of $D/r$ are its four rectangle corners, with signed infinite
limits at zero. Add $[a,b]$, retain both sign components, and merge only
overlapping closed intervals. If $0\in R$ and $0\in I_D$, the rank-one
equation permits every real $v$, so return $\mathbb R$. If $R=\{0\}$
and $0\notin I_D$, return the empty set: that primitive region is
incompatible with the model. Weighted interval sums then produce
Equation~\eqref{s61:projection}. This finite interval arithmetic preserves
unbounded rays and disconnected components.

For example, $S=1.5$, $D=.5$ and $p\in[.4,.6]$ give
\[
 \mathcal V_s=(-\infty,-1]\ \cup\ [4,\infty).
\]
Neither a finite interval joining plotted points nor a selected positive
branch describes this information. As an independent check on the algorithm,
fix a candidate $v$ and ask whether the rectangle image
$\{r(v-S):r\in R,S\in[a,b]\}$ intersects $[c,d]$.
Its endpoints are four products, so this test uses no ratio. The executed
16,800 comparisons agree, including exactly singular regions, and a finite
enumeration over all binomial counts for sizes 1--25 and 201 probability
values gives minimum audit coverage .95096 for nominal .95. These checks
validate finite identities; they are not a proof of the calendar CLT.

\subsection{Precision, alternatives and the bank's next decision}
A second valid projection uses separate bands for $(m_{s1},m_{s2})$ and
the original inverse map. It ignores much of their correlation when forming
a rectangle. The sum/difference construction uses that covariance before
forming bands: a common monetary shock cancels from the difference, while
uncertainty in the sum need not be amplified by $1/|2p-1|$.
Neither rectangular construction is uniformly shorter for every covariance
law. The registered comparison therefore evaluates both on the same draws,
alongside the regular delta interval and its availability rule.

A joint ellipsoidal response region could exploit further covariance but
requires reliable constrained projection, including unbounded directions.
Alternatively, for proposed $(v,p)$ one can test
$\widehat D-(2p-1)(v-\widehat S)$ using its candidate-specific covariance
and combine the audit evidence before eliminating $p$. Its variance includes
$V_{DD}+(2p-1)^2V_{SS}+2(2p-1)V_{SD}$; using a fixed variance at a fitted
inverse would recreate the weak-normalization error. A justified joint
test could improve power, but its critical value, nuisance elimination and
degenerate cases require a separate derivation. The selected rectangle
projection buys transparent coverage and exact topology at the cost of
conservatism. A joint Bayesian linkage/outcome model
\citep{guha2022} instead uses a different likelihood, feedback and prior
restriction; its posterior is not this confidence set.

The five-regime funding comparison uses 720 evaluation dates, an independent
2,880-date nuisance pilot and 200 independent audits per stratum, with 400
repetitions in each regime. Its true financial value is about 375,149 currency
units. The response law, target and error allocation remain fixed while the
mapping probabilities change.

\begin{table}[htbp]\centering\small
\caption{Weak linkage and financial inference. Projection width is the median
conditional on a bounded set, in millions of currency units. Regular coverage
counts unavailable intervals as no coverage.}
\label{s61:projection-results}
\begin{tabular}{lrrrrr}\toprule
$(p_0,p_1)$ & Regular & Regular & Projection & Bounded & Projection\\
 & available & coverage & coverage & fraction & width\\\midrule
(0.95, 0.80) & 400/400 & 0.9250 & 1.000 & 1.0000 & 1.498 \\
(0.70, 0.65) & 397/400 & 0.9150 & 1.000 & 0.9475 & 2.947 \\
(0.55, 0.52) & 15/400 & 0.0350 & 1.000 & 0.0000 & -- \\
(0.50, 0.50) & 2/400 & 0.0025 & 1.000 & 0.0000 & -- \\
(0.20, 0.35) & 395/400 & 0.9225 & 1.000 & 0.9550 & 2.542 \\
\bottomrule\end{tabular}
\end{table}

Projection covers the retained truth in all 400 draws in every regime; a
95\% binomial Monte Carlo interval for each coverage rate is $[.9908,1]$.
That result is conservative, not evidence of exact 100\% coverage. The weak
and singular cases return the whole real line in every draw. The moderate
case contains one disconnected result and twenty whole-line results; the
reversed case has eighteen whole-line results. Conditional median widths
for the candidate-mean rectangle are 1.910, 6.644 and 5.197 million in the
strong, moderate and reversed regimes, versus 1.498, 2.947 and 2.542 million
for sum/difference projection. Their conditional comparison excludes
unbounded sets; it cannot be interpreted as an unconditional width guarantee.

The regular procedure is substantially narrower where available, but it
covers .925 of the strong-regime draws, .915 of all moderate requests and
.9225 of all reversed requests. In the last two regimes its coverage
conditional on availability is .9219 and .9342. Even the strong-regime
finding is an adverse finite-sample result alongside the earlier
100/400-audit study; a new independent seed set does not have to reproduce
its favorable coverage. In the weak case the regular rule produces only
15 intervals, 14 covering truth; in the singular case it produces two,
only one covering truth. Those conditional denominators cannot justify
promoting the screened interval to inference robust to weak linkage.


An unbounded set can be the correct answer to the supplied information.
If an operationally justified lower bound on mapping accuracy excludes
$1/2$, it changes the admissible model and may recover a bounded result.
More representative correctness-count audits can resolve a genuinely nonzero
$r_s$, but cannot invert these mean equations when $p_s$ is exactly $1/2$.
In that case stronger identifiers, an additional measurement, or a changed
economic target such as the portfolio sum is required. Mislabelled audits,
selected audits and a persistent crosswalk outside the declared per-origin
map law require different observation models. High coverage obtained by
returning $\mathbb R$ is honest lack of information, not a useful investment
case or evidence of universal causal inference.

\section{Persistent customer identities and the price of new information}
\label{c7:persistent-audits}
A correctness count is not the same observation as an audited customer
history. Suppose a deposit team sees two candidate balance histories for
each customer pair. A fixed alias can swap these histories for every month.
Adding months then adds outcome information but no independent evidence about
which customer is which. If the probability of the correct map is one half,
the preceding mean equations identify the pair's total response and leave
its allocation unrestricted. Neither an inverse regularization nor a larger
sample of correctness counts changes that algebra.

Consider instead a declared population of independent pairs, indexed by
$i=1,\ldots,n$, with a fixed follow-up length $H$. The focal customer has
randomized actions $A_{it}\in\{-1,1\}$, independent of the potential outcome
disturbances, and a persistent map $M_i\in\{0,1\}$. Define
\[
 U_i=H^{-1}\sum_{t=1}^H A_{it}Y_{i,\mathrm{focal},t},\qquad
 \theta=\E U_i.
\]
For $Y_{i,\mathrm{focal},t}=b_i A_{it}+e_{it}$ with the stated randomization,
$\theta=\E b_i$ is the mean unit-action response. It is half the response
contrast between actions $+1$ and $-1$; this normalization must accompany
the monetary account. Within-pair random slopes, persistent disturbances and
repeated maps are allowed. Common shocks across different pairs are excluded
from this sampling experiment, not made harmless by the word ``cluster.''

Let $R_i$ indicate audit inclusion. An audit reveals the actual map and the
linked focal outcomes, with known inclusion probability $q_i>0$. Conditional
on the complete pair, $\Pr(R_i=1)=q_i$; inclusion may depend on declared
pre-audit features if its probability is known. Let $T_i$ be the corresponding
action--outcome moment using the mean of the two candidate histories. Then
\begin{equation}
 Z_i=T_i+\frac{R_i}{q_i}(U_i-T_i),\qquad
 \E[Z_i\mid\text{complete pair}]=U_i.
 \label{c7:audit-score}
\end{equation}
Thus $\widehat\theta=\bar Z$ is unbiased for the stated population moment.
If pairs are identically distributed, $q_i\ge q_{\min}>0$ and
$\E Z_i^2<\infty$, the ordinary independent-pair CLT gives
\[
 \widehat{\operatorname{se}}(\widehat\theta)^2
   =\frac{1}{n(n-1)}\sum_i(Z_i-\bar Z)^2.
\]
There are $n$ contributions, not $nH$ independently audited identities.
The augmentation $T_i$ is fixed by the candidate histories and does not
require a learned nuisance fit. Pure inverse-audit weighting uses $T_i=0$;
its expectation is the same but its variance can differ. An audited-only
mean generally changes the population under unequal inclusion.

The audit discussion of \citet[Section 2.1]{slawski2025} assumes independent
membership and accurate match labels. Equation~\eqref{c7:audit-score} is a
project construction for known unequal inclusion and persistent pairs; it
does not import that source's linked-data assumptions unchanged. A joint
linkage/outcome likelihood, such as the substantive alternative in the
preceding discussion, uses more distributional structure. Its possible
efficiency gain must be weighed against that additional model, rather than
treated as an automatic refinement of this moment argument.

\subsection{Label error is another observation law}
Suppose an audited label $L_i$ flips $M_i$ with a known independent symmetric
probability $e<1/2$. If $V_{i0},V_{i1}$ are the two candidate moments, then
\[
 \widetilde U_i=\frac{V_{i0}+V_{i1}}2+
       \frac{2L_i-1}{1-2e}\frac{V_{i0}-V_{i1}}2
\]
has conditional expectation $U_i$. Replacing $U_i$ by $\widetilde U_i$ in
Equation~\eqref{c7:audit-score} preserves unbiasedness and increases noise.
The denominator illustrates the information cost as $e$ approaches one half.
The error rate must be identified by a separate gold-standard design. Two
uncertain assertions do not identify it. If $e$ is estimated, its influence
and dependence on the audit sample must be included; the executed reference
treats it as externally known. An unavailable error law produces a request
for information, not an internally estimated truth label.

\subsection{A reproducible deposit decision}
The synthetic experiment has 800 independent pairs, eight dates, focal
population response $1.2$ and companion response $-0.4$. A pair's focal slope
also depends on a balanced binary feature and a random slope; its map is fixed
and has correctness probability one half. Audit inclusion is
$q_i=q(1+0.5X_i)$ for $X_i\in\{-1,1\}$. The nominal $q$ values are
$0.1,0.3,0.6$. A unit response is worth 1,000 currency units in the declared
fixed account, so the synthetic focal truth is 1,200. It is not a measured
deposit return for any bank.

Across 2,000 independent repetitions in each design, the augmented intervals
cover in 1,895, 1,895 and 1,894 draws respectively. Their mean monetary widths
are 413.27, 247.81 and 184.36. The corresponding 95\% Monte Carlo interval
for the first coverage frequency is $[.9368,.9569]$. With 15\% symmetric
label errors and the correct externally supplied error rate, coverage is
$1905/2000=.9525$ and mean width is 344.21. Falsely supplying the original
inclusion law when actual inclusion also depends on the feature reduces
coverage to $82/2000=.041$. This is evidence against ignoring audit selection,
not a limitation repaired by increasing the number of dates.

The richer audit design therefore supplies useful information in a case where
correctness counts and mean mixing alone do not. It does so by observing the
missing map and outcomes, not by extracting an identified contrast from the
old singular equations. Many-to-many ownership still needs an allocation
target, and an unknown selected audit still needs a sampling model. A source
owner's false semantic assertion can remain observationally undetectable.
The operation preserves that boundary and rejects a proposed combination
with common-calendar funding uncertainty until a joint sampling argument
has been supplied.

\section{Using covariance without hiding singular identification}
\label{c7:joint-projection}
The rectangular confidence construction is transparent but discards response
covariance when forming the financial set. A joint region can exploit that
covariance while retaining the same identification restrictions. Write the
four monetary response coordinates as
$x=(s_0,d_0,s_1,d_1)'$, with candidate average $s_j$ and half-difference $d_j$
in stratum $j$. Let $m$ estimate $x$ and let $V$ estimate its covariance.
Assume a fixed full-rank four-dimensional response limit and consistent $V$.
Allocate half of error probability $\alpha$ to the ellipsoid
\[
 \mathcal E=\{x:(x-m)'V^{-1}(x-m)\le
                    k\},\qquad k=\chi^2_{4,1-\alpha/2},
\]
and error $\alpha/4$ to each of the two exact binomial audit intervals.
Finite-sample exactness of the audit region does not turn a fitted calendar
response ellipsoid into an exact finite-sample pivot.

At fixed nonzero ranks $r_j=2p_j-1$, the focal financial contribution is
$c(r)'x-K$, where $K$ is the fixed direct cost and
\[
 c(r)=\left(\frac12,\frac1{2r_0},\frac12,\frac1{2r_1}\right)' .
\]
The ellipsoid's lower and upper support values are
\begin{equation}
 L(r)=c(r)'m-\sqrt{k\,c(r)'Vc(r)}-K,
 \quad U(r)=c(r)'m+\sqrt{k\,c(r)'Vc(r)}-K.
 \label{c7:ellipsoid-support}
\end{equation}
This follows by writing $x=m+V^{1/2}z$ and maximizing a linear functional
over $\|z\|\le\sqrt{k}$. It is a project derivation for the selected
response region, not a claim that one region shape uniformly dominates another.

\subsection{A global calculation over audit uncertainty}
If neither audit interval contains one half, each reciprocal rank $1/r_j$
ranges over a compact interval. Their product is a rectangle. The upper
support function $U$ is convex in these reciprocal coordinates because it is
the sum of a linear function and a norm; $L$ is concave. Every point in a
rectangle is a convex combination of its vertices. Convexity gives
$U(u)\le\max_v U(v)$ and concavity gives $L(u)\ge\min_v L(v)$.
Consequently the four corners give the global enclosing interval. No local
optimizer or finite effect grid can silently remove a tail.

When an audit interval touches one half, the implementation conservatively
retains the whole line. Some such cases admit a sharper disconnected result;
the original rectangle routine retains those rays and remains a challenger.
The ellipsoid implementation does not claim sharpness in this case. A singular
empirical response covariance also receives a declared conservative fallback,
rather than a pseudoinverse labelled as a four-dimensional confidence pivot.
Special-function inversions and endpoint rounding use ordinary floating-point
arithmetic. The global geometric argument is exact; a certified bound on
floating-point error has not been supplied.

If the true response is in $\mathcal E$ and both true audit probabilities
are in their intervals, the true financial contribution belongs to the
projected set. The union bound therefore gives asymptotic coverage at least
$1-\alpha$ under the stated response approximation, without a lower bound
on linkage rank. The set can be valid and financially uninformative.

\subsection{What the comparisons diagnose}
The preserved 400-draw strong-calendar design permits a retrospective
comparison because each draw retains all four means and their covariance.
The median bounded width falls from 1,498,486.81 for the rectangle to
1,202,977.31 for the ellipsoid; coverage is $400/400$ and $399/400$.
Moderate and reversed designs also have smaller median bounded widths in
this comparison. Every retained weak and singular ellipsoid remains the
whole line. These observations do not establish uniform dominance or
confirmatory calibration on newly selected calendar data.

The regular strong-design interval covers $370/400=.925$. Supplying the true
audit probabilities gives $372/400=.930$, and using the oracle first-order
linearization gives $371/400=.9275$. Thus estimated audit probabilities alone
do not explain the failure; the response approximation and its finite sample
behavior remain important. These arms identify contributors, not a unique
causal diagnosis of the coverage error. With freshly simulated Gaussian
primitive responses of known covariance, strong-design regular coverage is
$1920/2000=.960$, while both conservative projection constructions cover
all 2,000 draws. That primitive experiment isolates geometry from the original
calendar covariance estimation; it cannot retrospectively validate the latter.

The practical alternatives are consequently clear. Use the original rectangle
when preserving disconnected topology matters; use the registered ellipsoid
when covariance is well qualified and its outer interval is suitable; retain
the regular calculation as a disclosed research route; or purchase the
outcome-linked information in Section~\ref{c7:persistent-audits}. None may be
selected after observing which produces the narrowest favorable answer.

